\documentclass[letterpaper,11pt]{article}

\usepackage{amsmath, amsfonts, amssymb, amsthm}
\usepackage{lmodern}
\usepackage{enumitem}
\usepackage{hyperref}
\usepackage[x11names, rgb]{xcolor}
\usepackage{graphicx}
\usepackage{fullpage}
\usepackage{etoolbox}
\usepackage{mathtools}


\newcommand{\cX}{\mathcal{X}}
\newcommand{\cC}{\mathcal{C}}
\newcommand{\cF}{\mathcal{F}}
\newcommand{\R}{\mathbb{R}}
\newcommand{\dd}{{\rm d}}
\DeclareMathOperator{\sech}{sech}


\setlength{\parindent}{0pt}
\setlength{\parskip}{6pt}

\newcommand{\nopagenumbers}{
    \pagestyle{empty}
}

\newenvironment{indented}{
    \begin{list}{}{\leftmargin=1em}
    \item[]
}{
    \end{list}
}

\newtoggle{questionnumbers}
\toggletrue{questionnumbers}

\newcounter{questionCounter}
\newcounter{partCounter}[questionCounter]

\newcommand{\question}[1]{%
    \stepcounter{questionCounter}%
    \setcounter{partCounter}{0}%
    \vspace{0.25in}
    \hrule
    \vspace{0.4em}
    \noindent\textbf{Question \thequestionCounter:\ }#1
    \vspace{0.8em}
    \hrule
    \vspace{0.1in}
}

\newcommand{\qpart}{%
    \stepcounter{partCounter}%
    \begin{indented}
    \textbf{(\alph{partCounter}) }
}

\newcommand{\closepart}{
    \end{indented}
}

\newcommand{\myhwname}{AMATH 403 HW 7}
\newcommand{\myname}{Jordan Tucker}

\newcommand{\header}{
    \begin{center}
    {\Large \myhwname}\\
    \myname\\
    \today
    \end{center}
}

\nopagenumbers
% Document
\begin{document}
    \header
    \question{8.1.4}

    \qpart{Use the Fourier transform to solve the initial value problem
        \[
            \frac{\partial u}{\partial t}
            =
            \frac{\partial^2 u}{\partial x^2},
            \qquad
            u(0,x)=\delta'(x-\xi),
            \qquad
            -\infty < x < \infty,
            \quad
            t>0,
        \]
        whose initial data is the derivative of the delta function at a fixed
        position \(\xi\).
    }
    \begin{gather*}
        u_t = u_{xx}\\
        \cF{u_t} = \frac{\partial \hat{u}}{\partial t}\\
        \cF{u_{xx}} = (ik)^2 \hat{u} = -k^2\hat{u}
    \end{gather*}
    Form and solve the ODE
    \begin{align*}
        \frac{\partial \hat{u}}{\partial t} &= -k^2 \hat{u}\\
        \frac{\partial \hat{u}}{\hat{u}} &= -k^2 \partial t\\
        \ln(\hat{u}) &= -k^2 t + A(K)\\
        \hat{u}(t,k) &= A(k)e^{-k^2 t}
    \end{align*}
    Transform initial condition
    \begin{gather*}
        \hat{f}(k) = \frac{1}{\sqrt{2\pi}} \int_{-\infty}^{\infty} f(x)e^{-ikx} \, dx\\
        \hat{u}(0,k) = \frac{1}{\sqrt{2\pi}} \int_{-\infty}^{\infty}  \delta'(x-\xi) e^{-ikx} \, dx\\
    \end{gather*}\
    \[
        f(x) = e^{-ikx},\ f'(x) = -ike^{-ikx}
    \]
    \begin{gather*}
        \hat{u}(0,k) = \frac{1}{\sqrt{2 \pi}} \left(- (-ike^{-ik \xi}) \right) = \frac{ik}{\sqrt{2 \pi}} e^{-ik\xi}\\
        \hat{u}(0,k) = A(k)e^0 = A(k) \\
        A(k)= \frac{ik}{\sqrt{2\pi}}e^{-ik\xi}\\
        \hat{u}(t,k) = \frac{ik}{\sqrt{2\pi}}e^{-ik\xi} e^{-k^{2}t}
    \end{gather*}
    Use inverse fourier transform
    \begin{gather*}
        \hat{u}(t,k) =ik \left(\frac{1}{\sqrt{2\pi}}e^{-ik\xi} e^{-k^{2}t} \right)\\
        \hat{v}(t,k) =\frac{1}{\sqrt{2\pi}}e^{-ik\xi} e^{-k^{2}t}\\
        v(t,k) = \cF^{-1} \left(\frac{1}{\sqrt{2\pi}}e^{-ik\xi} e^{-k^{2}t} \right) = \frac{1}{\sqrt{4 \pi t}}e^{- \frac{(x-\xi)^2}{4t}}\\
    \end{gather*}
    Since $ik$ means take the spacial derivative:
    \begin{gather*}
        u(t,x) = \frac{\partial}{\partial x} v(t,x)\\
        u(t,x) = \frac{\partial}{\partial x} \left(\frac{1}{\sqrt{4 \pi t}}e^{- \frac{(x-\xi)^2}{4t}}\right)\\
        u(t,x) = \frac{1}{4\pi t} \left( e^{- \frac{(x-\xi)^2}{4t}}\right)\left(- \frac{x-\xi}{2t}\right)\\
        u(t,x) = -\frac{x-\xi}{4\sqrt{\pi} t^{3/2}} e^{-\frac{(x-\xi)^2}{4t}}
    \end{gather*}
    \closepart

    \qpart{
        Show that your solution can be written as the derivative
        \(\partial F / \partial x\) of the fundamental solution
        \(F(t,x;\xi)\).
        Explain why this observation should be valid.}
     \begin{gather*}
         F(t,x;\xi) = \frac{1}{\sqrt{4\pi t}} e^{- \frac{(x-\xi)^2}{4t}}\\
         u(t,x) = -\frac{x-\xi}{4\sqrt{\pi} t^{3/2}} e^{-\frac{(x-\xi)^2}{4t}}
     \end{gather*}
    Take the spacial derivative of $F$
    \begin{gather*}
        \frac{\partial F}{\partial x} = \frac{1}{\sqrt{4\pi t}} \left(- \frac{x-\xi}{2t} \right) \left(e^{-\frac{(x-\xi)^2}{4t}}\right)
        = -\frac{x-\xi}{4\sqrt{\pi} t^{3/2}} e^{-\frac{(x-\xi)^2}{4t}}\\
        u(t,x) - \frac{\partial F}{\partial x} = \left(-\frac{x-\xi}{4\sqrt{\pi} t^{3/2}} e^{-\frac{(x-\xi)^2}{4t}}\right)
        - \left(-\frac{x-\xi}{4\sqrt{\pi} t^{3/2}} e^{-\frac{(x-\xi)^2}{4t}}\right) = 0
    \end{gather*}
    $\frac{\partial F}{\partial x}$ is a solution to the heat equation, and it fits the initial condition $u(0,x) = \delta'(x-\xi)$ because the
    heat equation is linear with constant coefficients.
    Because solutions to the initial value problem for the heat equation are unique, $\frac{\partial F}{\partial x}$ must be the exact same function as our solution $u(t,x)$.
    \closepart
    \newpage
    \question{(8.1.17) Find the fundamental solution for the cable equation $vt = \gamma v_{xx} - \alpha v$ on the real line.}
    Fundamental solution means
    \[
        v(0,x)  = \delta(x)\\
    \]
    We can make the guess $v(t,x) = e^{-\alpha t} u(t,x)$
    \begin{align*}
        v_t &= (-\alpha e^{-\alpha t} u) + e^{-\alpha t} (u_t)\\
        v_t &= e^{-\alpha t} (u_t - \alpha u)
    \end{align*}
    \begin{align*}
        v_x &= e^{-\alpha t} u_x\\
        v_{xx} &= e^{-\alpha t} u_{xx}
    \end{align*}
    Substitute back into ODE
    \begin{align*}
        e^{-\alpha t} (u_t - \alpha u) &= \gamma(e^{-\alpha t} u_{xx}) - \alpha (e^{-\alpha t} u(t,x))\\
        u_t - \alpha u &= \gamma u_{xx} - \alpha u\\
        u_t &= \gamma u_{xx}
    \end{align*}
    It looks just like the heat equation
    \begin{align*}
        v(0,x) &= u(0,x)\\
        \delta(x) &= u (0,x)
    \end{align*}
    $u(t,x)$ is the fundamental solution to the heat equation.
    \begin{gather*}
        u(t,x) = \frac{1}{\sqrt{4 \pi \gamma} t} e^{-\frac{x^2}{4 \gamma t}}\\
        v(t,x) = e^{-\alpha t} \left( \frac{1}{\sqrt{4 \pi \gamma t}} e^{-\frac{x^2}{4 \gamma t}} \right)\\
        v(t,x) = \frac{1}{\sqrt{4 \pi \gamma t}} e^{-\left( \frac{x^2}{4 \gamma t} + \alpha t \right)}
    \end{gather*}
    \newpage
    \question{(8.3.5) Consider the parabolic equation
        \[
            \frac{\partial u}{\partial t}
            =
            x \frac{\partial^2 u}{\partial x^2}
            +
            \frac{\partial u}{\partial x},
            \qquad
            1 < x < 2,
        \]
        with initial and boundary conditions
        \[
            u(0,x)=f(x),
            \qquad
            u(t,1)=\alpha(t),
            \qquad
            u(t,2)=\beta(t).
        \]
    }

    \qpart{
        State and prove a version of the Maximum Principle for this problem.
    }

    If $u(t,x)$ is a smooth solution to $u_t = xu_{xx} + u_x$ on the domain $R = \left\{(t,x):0 \leq t \leq T, \ 1\leq x \leq 2 \right\}$
    then the maximum value of $u$ is at a point on the boundary of $R$.

    We will prove this by contradiction.
    We can start by assuming $u(t,x)$ has its maximum value at some interior point $u(t_0,x_0)$.
    Since it is the maximum value, the derivatives must be equal to zero.
    \begin{gather*}
        u_t(t_0,x_0) = 0\\
        u_x(t_0,x_0) = 0
    \end{gather*}
    And the second derivative must be less than or equal to zero.
    \[
        u_{xx}(t_0,x_0) \leq 0
    \]
    We can plug in our values at $(t_0,x_0)$ in to our ODE:
    \begin{align*}
        u_t(t_0,x_0) &= x u_{xx}(t_0,x_0) + u_x(t_0,x_0)\\
        0 &= x (\leq 0) + 0
    \end{align*}
    This simplifies down to $0 \leq 0$, which is true but does not provide a full contradiction.

   We can define $v(t,x)$ which is perturbed:
    \[
        v(t,x) = u(t,x) - \epsilon t
    \]
    where $\epsilon > 0$ is an arbitrarily small number.
    Since $u(t,x) = v(t,x) + \epsilon t$, we can find the derivatives:
    \begin{gather*}
        u_t = v_t + \epsilon \\
        u_x = v_x \\
        u_{xx} = v_{xx}
    \end{gather*}
    Subbing these into the original PDE gives us the PDE for $v$:
    \[
        v_t + \epsilon = x v_{xx} + v_x \implies v_t = x v_{xx} + v_x - \epsilon
    \]
   Now we can assume $v(t,x)$ has its maximum at some interior point $(t_0, x_0)$.
   The spatial derivatives are zero and less than 0:
    \begin{gather*}
        v_x(t_0,x_0) = 0\\
        v_{xx}(t_0,x_0) \leq 0
    \end{gather*}
    For the time derivative, since we are defining our domain up to time $T$, $v_t$ could be zero positive if the maximum happens
    right at the very end of the time limit $T$.
    \[
        v_t(t_0,x_0) \geq 0
    \]
    We can substitute these values at $(t_0, x_0)$ into our new ODE for $v$:
    \begin{align*}
        v_t(t_0,x_0) &= x_0 v_{xx}(t_0,x_0) + v_x(t_0,x_0) - \epsilon\\
        v_t(t_0,x_0) &= x_0 (\leq 0) + 0 - \epsilon
    \end{align*}
    Since our domain is $1 \leq x \leq 2$, we know $x_0$ is a strictly positive number.
    A positive number multiplied by $\leq 0$ is still $\leq 0$. This leaves us with:
    \[
        v_t(t_0,x_0) \leq -\epsilon
    \]
    Since $\epsilon$ is a positive number, $-\epsilon$ is negative, which means $v_t(t_0,x_0)$ must be strictly less than zero.
    but we already established that $v_t(t_0,x_0) \geq 0$.

    We have $0 \leq v_t \leq -\epsilon < 0$, which is a contradiction.

    Because of this contradiction, our assumption is false.
    So $v(t,x)$ cannot have its maximum on the interior, and must have it on the boundary.
    If we let $\epsilon \to 0$, $v(t,x)$ turns back into $u(t,x)$, which means $u(t,x)$ must also have its maximum on the boundary.
    \closepart
    \qpart{
        Establish uniqueness of the solution to this initial-boundary value problem.
    }

    Let's assume we have two solutions, $u_1(t,x)$ and $u_2(t,x)$, that both solve the PDE and
    the given initial and boundary conditions.

    We can define a new functions $w(t,x)$ such that:
    \[
        w(t,x) = u_1(t,x) - u_2(t,x)
    \]

    The PDE is linear, so we can subtract the PDE for $u_2$ from the PDE for $u_1$, which
    shows us that $w(t,x)$ will also solve the PDE:
    \[
        w_t = x w_{xx} + w_x
    \]

    Now we can check the initial and boundary conditions for $w$.
    Since $u_1$ and $u_2$ share the exact
    same conditions, they will subtract to zero:
    \begin{gather*}
        w(0,x) = u_1(0,x) - u_2(0,x) = f(x) - f(x) = 0 \\
        w(t,1) = u_1(t,1) - u_2(t,1) = \alpha(t) - \alpha(t) = 0 \\
        w(t,2) = u_1(t,2) - u_2(t,2) = \beta(t) - \beta(t) = 0
    \end{gather*}
    This means that on the entire boundary,
    the value of $w(t,x)$ is $0$.

    Using the maximum principle, the maximum value of $w(t,x)$ must
    be on this boundary.
    So $w(t,x) \leq 0$ everywhere on the interior domain.

    We can also use the maximum principle on $-w(t,x)$
    to show that the minimum value of $w(t,x)$ must also be on the boundary.

    Since the boundary is 0, this means $w(t,x) \geq 0$ everywhere on the interior.

    Since $0 \leq w(t,x) \leq 0$, $w(t,x) = 0$ everywhere on our domain.

    Therefore, $u_1(t,x) - u_2(t,x) = 0 \implies u_1(t,x) = u_2(t,x)$ which shows that any two
    solutions to the PDE must be identical, so the solution is unique.
    \closepart
\end{document}